\section{Contributions}
\begin{itemize}
\item We provide the first complete implementation of a probabilistic separation logic.
Specifically Lilac is formalised in the Lean theorem prover complete with both soundness proof and
extensive support for tactic-based weakest-precondition style reasoning about programs
\item One of the first non-trivial instantiation of the Lean port of the Iris framework (originally developed in Rocq).
\item A Novel treatment of infinite looping programs which may construct a convergent sequence of values.
The correctness of some algorithms for probabilistic inference is given by convergence of such a sequence.
So add a connective to Lilac for reasoning about infinite loops and re-establish the soundness of the logic.

``The search for well-specified probabilistic models—models that correctly describe the pro-
cesses that produce data—is one of the enduring ideals of the statistical sciences. Yet, in
only the simplest of settings are we able to achieve this goal'' - Papamakarios et al

``Although our tactics perform symbolic execution in small steps, one can
easily define a tactic that performs symbolic execution repeatedly. However,
in concurrent separation logic, small step symbolic execution is often needed
so as to be able to open invariants around atomic expressions''
 - Krebbers et al. 2017b https://iris-project.org/pdfs/2017-popl-proofmode-final.pdf

 so it makes sense for be to automatically apply wp\_let, wp\_ret and wp\_unif and wp\_ber
 immediately. However is it not possible for this to get stuck if we are not dealing
 with unif and ber alone (eg. complicate (det/rand)val coming from the meta logic??)

Shall I talk about logic programming via type classes? Probably briefly should do. To
illustrate why and how certain type classes were implemented for certain propositions.
Especially interesting is this line [Krebbers et al. 2017b] 
``However, since we use a shallow embedding
to represent the propositions of our object logic, propositions are
semantic objects, so we cannot just write a Coq function to perform this traversal. Instead, we use logic programming using type classes
to perform such manipulations on the meta level''

This work is a good example exhibiting an easy to use separation logic proofmode, without
the complexity of step-indexing (as in most iris projects). This is because our language is first-order
and does not support recursion and yet has an orthogonal axis of intricacy which makes it interesting.
As such it makes a good case study for the application of modern proofmode construction
techniques decoupled from the step-indexed world of Iris in which practically every other
significant program logic lives.

Expressive tactic-based reasoning about programs is developed through Lean meta-programming. 
\end{itemize}

\section{Motivation}
\emph{Here we give an informal account of what probabilistic programs are and why we would care to use them}
A Probabilistic programming language is not an established or general paradigm of programming such as Functional or Imperative programming.
Most generally probabilistic programs are simply programs which have a \texttt{sample} construct for sampling from probability distributions,
and optionally, an \texttt{observe} for specifying a new distribution given an old distribution and some newly observed data.

Statisticians use languages such as Stan \cite{stan}, Hakaru \cite{hakaru} and Anglican \cite{anglican}, to build complex statistical models
succinctly using the expressivity of a programming language. Note that here, a model is simply a probability distribution over some quantity of
interest, take for instance the probability of a 3d structure being correct for the given protein sequence, (other examples may come from
weather forecasting, medical diagnosis, economics, statistical mechanics etc).
"Running" this code corresponds to running a simulation of the model, where we can get answers such as the mean or variance, or just samples
(sample protein structure to test in a lab in the example).

What distinguishes a probabilistic programming language is that it decouples the specification of a statistical model from the algorithm for
its simulation (henceforth called probabilistic inference). Inference algorithms are often highly optimised and specialised to the specific model
and the interpreter of a PPLs attempts to make the best choice/combination of inference algorithms for the given code (model), relieving
the user of this additional work.

So in other words PPLs are a layer of abstraction which sits between a statistical model and the inference algorithms which simulate it.
This abstraction unlocks the verification of correctness properties of the statistical model, which is the main topic of this report.

\subsection{Probabilistic Programs in the wild}
1) Cryptography -> we're not building a statistical model here... what's going on?
2)
\subsection{}

Importantly, probabilistic programming decouples the specification of a model from its calculation:
The specification is the source code, obtained by composing distributions \texttt{sample} and
\texttt{observe} operations. The calculation is from the inference algorithms which  

Statisticians use languages such as Stan \cite{stan}, Hakaru \cite{hakaru} and Anglican \cite{anglican}, to build complex statistical models
succinctly using the expressivity of a programming language.

this paradigm code represents statistical models which are use for simulation and prediction in complex systems  
Some examples of such languages include Stan \cite{stan}, Hakaru \cite{hakaru} and Anglican \cite{anglican} 